\subsection{Fortschritt-Visualisierung}
Um dem Nutzer Rückmeldung zu geben, wie weit der Algorithmus mit der Gipfelsuche fortgeschritten ist wird auf zwei seperate mechanismen gesetzt. Zum einen wird auf einen Ladebalken, wie er Beispielsweise im Windows Update Bildschirm zu sehen ist gesetzt. Zusätzlich werden ausführliche Logs in der Kommandozeile implementiert, die zudem auch für die Debugging Zwecke verwendet werden können. Die Kommandozeilenausgabe ist leit mit der Hilfe der Python Bibliothek logging zu implementieren. Dem Ladebalken hingegen muss immer mitgeteilt werden wie weit die Gipfel Analyse Fortgeschritten ist. Dies ist besonders schwierig, da zu Beginn die genaue Länge nicht abschätzbar ist und auch auf gleich großen Dateien und mit gleiochen eingabeparametern stark variert. Daher wird eine Klasse - Analysis State - implementiert, die von dem Algorithmus REückmeldung erhält, wie weit die Analyse fortgeschritten ist und dies in eine Prozentzahl übersetzt. Diese Übersetzung muss für Mehrdateiausführungen anders erfolgen wie für Einzeldateien. 
Im Grunde nutzt der Analysis State die Tatsache aus, dass ein Schritt immer die Anzahl der zu untersuchenden Erhebungen für den Folgeschritt bestimmt. Das bedeutet, dass die für die Prominenzberechnung zu überprüfenden Gipfel von der Lokalen Maxima Bestimmung entdeckt gefunden werden. Da vor der lokalen Maximabestimmung kein Schritt existiert, kannn jedoch hier keine Aussage über die zu erwartenden Werte getroffene werden (das ist an deiser stelle aber auch nicht sinvoll) \todo{erklären warum nicht}. \todo{anahme immer gleich lange ausführung pro element ist immer falsch}

- einzeldateien
- mehrfachdateien

- Fortschritte je Schritt
- Fortschritte 